\section{Concordancia básica}
\textbf{Número y género.}
\begin{itemize}[leftmargin=*]
  \item Plural: siempre \textit{-e}.\ \textit{De rode bloemen}.
  \item Neutro singular con artículo definido: \textit{het groene boek}.
  \item Neutro singular indefinido: sin \textit{-e} (excepto adjetivos fuertes): \textit{een groen boek}.
\end{itemize}

\textbf{Adjetivos fuertes comunes}
\begin{itemize}[leftmargin=*]
  \item \textit{mooi, groot, klein, nieuw, oud}. (Siguen la regla general de \textit{-e} salvo en neutro indefinido.)
\end{itemize}

\textbf{Mini práctica}
\begin{itemize}[leftmargin=*]
  \item Completa con \textit{-e} cuando corresponda en 6 frases.
  \item Pasa de singular a plural conservando la concordancia.
\end{itemize}
